\section{Define the resonance functional and its norm}
\label{sec:rigged-hilbert}

\subsection{Gaussian regularization of a Gamow state}
\label{subsec:zeldovich}
% BEGIN CURATED SUBJECT INDEX
\index{operator!adjoint and self-adjoint}
\index{boundary condition!incoming and outgoing}
\index{Gamow--Siegert state}
\index{Zel'dovich regularization}
\index{Gaussian regularization}
\index{Berggren completeness relation}
% END CURATED SUBJECT INDEX

Let $\psi_{\mathrm R}(x)$ obey the outgoing condition and grow exponentially
on the real axis.  Zel'dovich regularization defines a bilinear overlap by
inserting a Gaussian and analytically removing it:
\begin{equation}
  (\psi_{\mathrm R}\mathbin{|}\psi_{\mathrm R})_{\!Z}
  =\lim_{\eta\to0^+}
  \int_{-\infty}^{\infty}\rmd x\,
  \rme^{-\eta x^2}\psi_{\mathrm R}(x)^2 .
  \label{eq:zeldovich-norm}
\end{equation}
The square is not an absolute square.  It is the bilinear pairing appropriate
to a time-reversal-invariant resonance problem.  The integral is first
evaluated where $\eta>0$ makes it convergent and is then continued to the
desired parameter domain
\cite{Zel'dovich:1961,Berggren:1968zz}.

The regulator can be expressed as a nonunitary multiplication operator
\begin{equation}
  G_\eta=\rme^{-\frac{\eta x^2}{2}},
  \qquad
  H_\eta=G_\eta H G_\eta^{-1} .
  \label{eq:zeldovich-transform}
\end{equation}
For the kinetic operator, the transformed momentum is
\begin{equation}
  G_\eta pG_\eta^{-1}=p-\rmi\hbar\eta x,
  \qquad p=-\rmi\hbar\frac{\rmd}{\rmd x} .
  \label{eq:zeldovich-momentum}
\end{equation}
Consequently $H_\eta$ is non-self-adjoint.  The transformation and complex
scaling look different, but both move the asymptotic growth into the operator
and leave a regularized vector in a square-integrable representation.  Exact
WKB analysis makes their common analytic content particularly transparent
\cite{Morikawa:2025xjq}.

\subsection{Berggren's proof: orthogonality, normalization, and completeness}
\label{subsec:berggren-proof}
% BEGIN CURATED SUBJECT INDEX
\index{locally convex space}
\index{Jost function}
\index{partial-wave expansion}
\index{scattering!Coulomb}
\index{biorthogonality}
\index{Zel'dovich regularization}
\index{Berggren completeness relation}
% END CURATED SUBJECT INDEX

The Gaussian formula in Eq.~\eqref{eq:zeldovich-norm} becomes mathematically
useful only after three questions are answered: what analytic continuation is
being taken, why different resonant states are orthogonal in the resulting
bilinear pairing, and why the resonance sum must be accompanied by a deformed
continuum.  Berggren's construction answers these questions for radial
short-range problems
\cite{Berggren:1968zz,Michel:2009,MichelPloszajczak:2021}.  We give the
proof in enough detail to expose its assumptions.

\begin{theorem}[Finite-range weak Berggren expansion]
Fix a partial wave on $L^2(0,\infty)$ with the regular boundary condition at
the origin.  Let the real potential be locally integrable and vanish for
$r>R$.  Assume that its Jost function has no threshold zero, spectral
singularity, or multiple zero in the region swept between the positive real
$k$ axis and a piecewise smooth contour $L^+$ in the fourth quadrant.  Let
$\PhiSpace_{L^+}$ consist of radial test functions whose Jost transforms are
analytic in that region and decrease sufficiently that the closing arcs and
the small threshold arc vanish.

Then, as a weak identity from $\PhiSpace_{L^+}$ to its antidual,
\begin{equation}
  \sum_{n\in b\cup d}
  \ket{u_n}\bra{\widetilde u_n}
  +\int_{L^+}\rmd k\,
  \ket{u(k)}\bra{\widetilde u(k)}
  =\identity .
  \label{eq:berggren-theorem-statement}
\end{equation}
The set $b$ contains bound-state poles and $d$ contains precisely the simple
decaying poles between the two contours.  The discrete vectors are normalized
by the analytic continuation of the bilinear Zel'dovich pairing, and the
continuum uses the corresponding biorthogonal delta normalization.
\label{thm:berggren-finite-range}
\end{theorem}

The finite-range assumption lets the tail and the contour deformation be
checked directly.  Exponentially decreasing, Coulomb-modified,
coupled-channel, and nonlocal potentials have useful Berggren extensions, but
their analytic domains, threshold terms, and test spaces must be restated;
they are not automatic corollaries of
Theorem~\ref{thm:berggren-finite-range}.

\subsubsection*{The regulated outgoing integral}
% BEGIN CURATED SUBJECT INDEX
\index{boundary condition!incoming and outgoing}
\index{Stokes curve}
\index{Zel'dovich regularization}
% END CURATED SUBJECT INDEX

First consider the elementary tail integral
\begin{equation}
  I_\eta(q;R)
  =\int_R^\infty\rmd r\,
  \rme^{-\eta r^2+\rmi q r},
  \qquad \eta>0 .
  \label{eq:zeldovich-elementary-integral}
\end{equation}
Completing the square gives
\begin{equation}
  I_\eta(q;R)
  =\frac{\sqrt{\pi}}{2\sqrt{\eta}}
  \rme^{-\frac{q^2}{4\eta}}
  \operatorname{erfc}\left(
  \sqrt{\eta}R-\frac{\rmi q}{2\sqrt{\eta}}
  \right) .
  \label{eq:zeldovich-erfc}
\end{equation}
For every fixed $\eta>0$ this is an entire function of $q$.  In the upper
half-plane $\im q>0$, the regulator can be removed by ordinary dominated
convergence:
\begin{equation}
  \lim_{\eta\to0^+}I_\eta(q;R)
  =\int_R^\infty\rmd r\,\rme^{\rmi q r}
  =\frac{\rmi}{q}\rme^{\rmi qR} .
  \label{eq:zeldovich-tail-upper}
\end{equation}
The right-hand side is analytic for $q\ne0$.  Its unique analytic
continuation defines the Zel'dovich value in a resonance sector where the
unregulated integral diverges:
\begin{equation}
  \left[\int_R^\infty\rmd r\,\rme^{\rmi q r}\right]_{Z}
  :=\left[\frac{\rmi}{q}\rme^{\rmi qR}\right]_{\mathrm{a.c.}} .
  \label{eq:zeldovich-tail-definition}
\end{equation}
Writing ``take $\eta\to0$'' without the analytic-continuation qualification
is incomplete: different asymptotic sectors of the complementary error
function are separated by Stokes lines.  The prescription is fixed by first
working in the convergent half-plane and then continuing the resulting
analytic function.

\subsubsection*{Orthogonality on a finite-range $s$-wave model}
% BEGIN CURATED SUBJECT INDEX
\index{boundary condition!incoming and outgoing}
\index{Gamow--Siegert state}
\index{Riccati equation}
\index{complex-symmetric operator}
\index{Berggren completeness relation}
\index{exceptional point}
\index{self-orthogonality}
% END CURATED SUBJECT INDEX

Let $V(r)=0$ for $r>R$ and write the reduced $s$-wave equation as
\begin{equation}
  \left[-\frac{\rmd^2}{\rmd r^2}+U(r)\right]u_n(r)
  =k_n^2u_n(r),
  \qquad U(r)=\frac{2\mu}{\hbar^2}V(r),
  \label{eq:berggren-radial-equation}
\end{equation}
with $u_n(0)=0$ and the purely outgoing tail
\begin{equation}
  u_n(r)=u_n(R)\rme^{\rmi k_n(r-R)},
  \qquad r\geq R .
  \label{eq:berggren-outgoing-tail}
\end{equation}
For two distinct solutions, Green's identity gives
\begin{align}
  &(k_n^2-k_m^2)\int_0^R\rmd r\,u_n(r)u_m(r)
  \notag\\
  &\qquad=
  \left[u_nu_m'-u_mu_n'\right]_{r=R}
  =\rmi(k_m-k_n)u_n(R)u_m(R) .
  \label{eq:berggren-green-identity}
\end{align}
Consequently
\begin{equation}
  \int_0^R\rmd r\,u_nu_m
  =-\frac{\rmi}{k_n+k_m}u_n(R)u_m(R) .
  \label{eq:berggren-interior-overlap}
\end{equation}
On the other hand, Eq.~\eqref{eq:zeldovich-tail-definition} with
$q=k_n+k_m$ gives
\begin{equation}
  \left[\int_R^\infty\rmd r\,u_nu_m\right]_{Z}
  =\frac{\rmi}{k_n+k_m}u_n(R)u_m(R) .
  \label{eq:berggren-tail-overlap}
\end{equation}
The interior and exterior terms cancel exactly:
\begin{equation}
  (u_n\mathbin{|}u_m)_Z
  :=\int_0^R\rmd r\,u_nu_m
  +\left[\int_R^\infty\rmd r\,u_nu_m\right]_{Z}
  =0,
  \qquad k_n^2\ne k_m^2 .
  \label{eq:berggren-orthogonality}
\end{equation}
No complex conjugation appears.  The pairing is the $c$ product appropriate
to a time-reversal-invariant complex-symmetric radial problem.  Ordinary
Hermitian orthogonality would be incompatible with two outgoing states.

For $m=n$, the same prescription gives the finite Gamow norm
\begin{equation}
  \mathcal N_n^2
  =(u_n\mathbin{|}u_n)_Z
  =\int_0^R\rmd r\,u_n(r)^2
  +\frac{\rmi}{2k_n}u_n(R)^2 .
  \label{eq:berggren-s-wave-norm}
\end{equation}
The norm is generally complex before a normalization convention is imposed.
One may rescale $u_n$ so that $\mathcal N_n^2=1$.  At an exceptional point
the bilinear norm can vanish; this self-orthogonality signals a higher-order
pole and cannot be repaired by ordinary rescaling.

For $l>0$, the outgoing tail is a Riccati--Hankel function rather than a
single exponential.  The proof is unchanged in structure.  If
\begin{equation}
  L_l^{(+)}(k;R)
  :=\frac{\partial_r f_l(k,R)}{f_l(k,R)}
  \label{eq:outgoing-log-derivative}
\end{equation}
is the outgoing logarithmic derivative, the analytically continued exterior
part of the norm is
\begin{equation}
  \left[\int_R^\infty\rmd r\,u_n(r)^2\right]_Z
  =\frac{u_n(R)^2}{2k_n}
  \left.\frac{\partial L_l^{(+)}(k;R)}{\partial k}
  \right|_{k=k_n} .
  \label{eq:berggren-general-tail-norm}
\end{equation}
For $l=0$, $L_0^{(+)}=\rmi k$, and
Eq.~\eqref{eq:berggren-general-tail-norm} reduces to the surface term in
Eq.~\eqref{eq:berggren-s-wave-norm}.  Formula
\eqref{eq:berggren-general-tail-norm} follows either by analytically
continuing the finite Riccati--Hankel expansion or by taking the coincident
limit of Green's identity.  This boundary form is often more stable
numerically than integrating a Gaussian at very small regulator.

\subsubsection*{Connection with the Jost derivative and the pole residue}
% BEGIN CURATED SUBJECT INDEX
\index{resolvent}
\index{Jost function}
\index{boundary condition!incoming and outgoing}
\index{resonance!residue}
\index{radial equation}
\index{Zel'dovich regularization}
\index{Berggren completeness relation}
% END CURATED SUBJECT INDEX

At a simple zero $F_l(k_n)=0$, the regular and outgoing solutions are
proportional.  Differentiate the radial equation with respect to $k$, pair it
bilinearly with $u_n$, and integrate by parts.  The interior contribution and
the outgoing surface term combine into Eq.~\eqref{eq:berggren-general-tail-norm}.
The result relates the regularized norm to the derivative of the incoming
Jost coefficient:
\begin{equation}
  \mathcal N_n^2
  \propto F_l'(k_n),
  \label{eq:berggren-jost-derivative}
\end{equation}
where the proportionality factor depends only on the chosen normalizations of
$\varphi_l$ and $f_l$.  Substitution into
Eq.~\eqref{eq:radial-green-residue} gives
\begin{equation}
  \Res_{E=E_n}G_l^{(+)}(E;r,r')
  =\frac{u_n(r)u_n(r')}{\mathcal N_n^2}
  \label{eq:berggren-residue-normalization}
\end{equation}
in the book's convention $\Resolvent(E)=(E-H)^{-1}$.  Thus Zel'dovich normalization is
the coordinate-space form of unit residue normalization.  This is the
mathematical reason it is compatible with scattering theory.

\subsubsection*{Contour deformation and the Berggren completeness relation}
% BEGIN CURATED SUBJECT INDEX
\index{Hilbert space}
\index{locally convex space}
\index{Jost function}
\index{boundary condition!incoming and outgoing}
\index{resonance!residue}
\index{Zel'dovich regularization}
\index{Berggren completeness relation}
\index{rigged Hilbert space}
% END CURATED SUBJECT INDEX

Begin with the real-energy radial completeness relation, normalized so that
the continuum states satisfy a Dirac delta:
\begin{equation}
  \sum_{b}\,u_b(r)u_b(r')
  +\int_0^\infty\rmd k\,u(k,r)u(k,r')
  =\delta(r-r') .
  \label{eq:newton-radial-completeness}
\end{equation}
The sum contains physical bound states.  For test functions whose momentum
representatives have sufficient analyticity and decay, the continuum kernel
can be continued into the fourth quadrant.  Deform the positive real
$k$ axis to a contour $L^+$ that passes below selected decaying resonance
poles but avoids threshold and other singularities.  Cauchy's theorem gives
the residues of every crossed Jost zero.  With the residue normalization of
Eq.~\eqref{eq:berggren-residue-normalization}, the result is
\begin{equation}
  \sum_{b,d}\ket{u_n}\bra{\widetilde u_n}
  +\int_{L^+}\rmd k\,
  \ket{u(k)}\bra{\widetilde u(k)}
  =\identity .
  \label{eq:berggren-completeness-proof}
\end{equation}
Here $b$ labels bound states and $d$ the decaying poles between the real axis
and $L^+$.  The tilde denotes the bilinear dual, not ordinary complex
conjugation.  The large arc vanishes only on an appropriate test space; hence
Eq.~\eqref{eq:berggren-completeness-proof} is naturally a weak identity in a
rigged Hilbert space.  If $L^+$ is returned to the real axis, the resonance
residues are reabsorbed into the continuum integral and
Eq.~\eqref{eq:newton-radial-completeness} is recovered.

This proof explains four features that are sometimes conflated.  First,
Zel'dovich regularization defines the bilinear norm of an outgoing residue.
Second, Green's identity proves orthogonality after the regulated tail is
included.  Third, Jost analyticity identifies the same norm with a pole
residue.  Fourth, contour deformation proves completeness only for the
combined set of bound states, enclosed resonances, and the $L^+$ continuum.
Dropping the continuum destroys the identity and removes threshold and
background physics.

\subsection{The Gelfand triple}
\label{subsec:gelfand-triple}
% BEGIN CURATED SUBJECT INDEX
\index{operator!adjoint and self-adjoint}
\index{Hilbert space}
\index{locally convex space}
\index{topology}
\index{spectrum!point}
\index{resolvent}
\index{direct integral}
\index{resonance!residue}
\index{Gamow--Siegert state}
\index{rigged Hilbert space}
% END CURATED SUBJECT INDEX

A rigged Hilbert space, or Gelfand triple, is a nested sequence
\begin{equation}
  \PhiSpace\subset\Hil\subset\PhiSpace^{\times} .
  \label{eq:gelfand-triple}
\end{equation}
The space $\PhiSpace$ is dense in $\Hil$ and carries a topology stronger than
the Hilbert norm.  The antidual $\PhiSpace^{\times}$ consists of continuous
antilinear functionals on $\PhiSpace$.  Generalized continuum eigenvectors and
Gamow vectors live in $\PhiSpace^{\times}$ rather than in $\Hil$.  The choice
of $\PhiSpace$ is part of the physics: it specifies localization, smoothness,
analyticity, and allowed contour deformations
\cite{DelaMadrid:2001wwa,DelaMadrid:2001dln,DelaMadrid:2002cz,Rafael:2012}.

The arrows in Eq.~\eqref{eq:gelfand-triple} are continuous embeddings, not
only set inclusions.  A Hilbert vector $\psi$ is embedded in the antidual by
$\phi\mapsto\dualpair{\phi}{\psi}$.  In the energy direct integral of
Eq.~\eqref{eq:spectral-direct-integral}, evaluation of an $L^2$ section at one
energy is not continuous and is not even defined independently of its
almost-everywhere representative.  A test topology imposing energy
regularity makes the evaluation functional continuous.  This elementary fact
is the rigorous content behind the notation $\braket{E,\alpha|\phi}$.
Analytic continuation of that evaluation functional can then define a Gamow
vector at complex energy, provided the chosen seminorms control the required
continuation.

For a test vector $\phi\in\PhiSpace$, a resonance functional can be defined by
the pole residue
\begin{equation}
  \dualpair{\phi}{\Psi_{\mathrm R}}
  =\Res_{z=z_{\mathrm R}}
  \dualpair{\phi}{\Resolvent(z)\chi} ,
  \label{eq:gamow-functional}
\end{equation}
where $\chi\in\PhiSpace$ is chosen so that the residue is nonzero.  The
functional satisfies a generalized eigenvalue equation
\begin{equation}
  \dualpair{H\phi}{\Psi_{\mathrm R}}
  =z_{\mathrm R}\dualpair{\phi}{\Psi_{\mathrm R}} .
  \label{eq:generalized-eigenvalue}
\end{equation}
No contradiction with self-adjointness arises because
$\Psi_{\mathrm R}\notin\Hil$ and Eq.~\eqref{eq:generalized-eigenvalue} is an
identity in the dual pairing, not an ordinary Hilbert-space eigenvalue
equation.

\subsection{Hardy classes and time direction}
\label{subsec:hardy-time}
% BEGIN CURATED SUBJECT INDEX
\index{locally convex space}
\index{Gamow--Siegert state}
\index{Hardy class}
% END CURATED SUBJECT INDEX

One can choose energy wave functions in Hardy classes analytic in an upper or
lower half-plane.  Contour closure then produces a semigroup time evolution
for Gamow vectors: decaying poles contribute for $t>0$, and growing
time-reversed poles contribute for $t<0$.  This construction gives a clean
mathematical encoding of preparation and detection, but it is not forced by
the Hilbert-space Hamiltonian alone.  It depends on the test space and on which
boundary values are identified with prepared states and observables.  We
therefore use the rigged-space formalism to define pole functionals without
claiming that a microscopic arrow of time follows solely from analytic
continuation.

\subsection{Relation to complex scaling}
\label{subsec:rhs-csm}
% BEGIN CURATED SUBJECT INDEX
\index{spectrum!point}
\index{resolvent}
\index{resonance!residue}
\index{Gamow--Siegert state}
\index{complex scaling}
% END CURATED SUBJECT INDEX

For analytic test vectors, the complex-scaled eigenvector and the Gamow
functional represent the same resolvent residue.  Schematically,
\begin{equation}
  \dualpair{\phi}{\Psi_{\mathrm R}}
  =\braket{U(\rmi\theta)\phi|
  \Psi_{\mathrm R}^{\theta}}_{\theta},
  \label{eq:rhs-csm-pairing}
\end{equation}
where $\Psi_{\mathrm R}^{\theta}\in L^2$ is the eigenvector of $H_\theta$ and
the right-hand side uses the left--right pairing appropriate to the deformed
operator.  Equation~\eqref{eq:rhs-csm-pairing} is independent of $\theta$ as
long as the pole remains exposed and the test vector is analytic under the
deformation.  The equality has been demonstrated explicitly in solvable
models and in the Friedrichs model
\cite{Antoniou:2001,Morikawa:2025xjq}.

This relation also explains why the number of discrete resonance eigenvalues
visible in $H_\theta$ depends on $\theta$ while the underlying pole set does
not.  Changing $\theta$ changes which poles are represented by $L^2$
eigenvectors and which remain behind the rotated continuum.  The rigged-space
functional exists through continuation beyond that particular representation,
subject to the analytic domain of the problem.

\subsection{Probability, observables, and normalization}
\label{subsec:resonance-observables}
% BEGIN CURATED SUBJECT INDEX
\index{wave packet}
\index{resonance!residue}
\index{Gamow--Siegert state}
\index{biorthogonality}
\index{exceptional point}
% END CURATED SUBJECT INDEX

An expectation value formed from a Gamow state can be complex.  This is not a
probability expectation in a closed system.  It is a pole residue or a
biorthogonal matrix element that may encode both a central value and a decay
scale.  Observable probabilities are obtained from real-energy scattering
states, wave packets, detector models, or a unitary dilation of the open
dynamics.  Resonance vectors are efficient intermediate objects in those
calculations.

Three normalizations should therefore be kept distinct.  Hilbert normalization
uses $\int|\psi|^2=1$ for a bound state or wave packet.  Continuum
normalization uses a Dirac delta, such as
$\braket{k|k'}=\delta(k-k')$.  Resonance normalization uses an analytic
bilinear prescription, a residue condition, or a biorthogonal pairing.  The
  singular transition between the last two normalizations is precisely what
  must be controlled when a resonance reaches the complex-scaled continuum at
  an exceptional point.
